% ============================================================
% M3 S2 导学件·波1 片件 —— 课时19 章末总结与复习（母版照抄制·补位臂4b）
% 生成：M3 成卷轮 S2 波1 补位臂4b｜2026-09-14｜编译 cwd＝本目录（中文路径禁 \input 绝对路径）
%   xelatex main-true.tex ×2 → 含详解印本；xelatex main-false.tex ×2 → 纯题版（单源双本）
% 输入链（全只读）：题面库 课时19-章末总结与复习.md（题面逐字装配）＋ -答案侧.md（值逐字回嵌）
%   ＋ 定稿/批E-课时19-章末总结与复习.md（详解回嵌源＝§七 补产 16 块【详解】＋件末「导学 G5 补产」5 键）；
%   toolchain＝qp-m3.sty（钉后版 md5 7c3930362be8a0a2bdf21bbf8ac16573，本地挂载副本，破损即停）。
% 母版体例：详见 成卷/导学件/课时01-坐标法/母版体例说明.md（骨架/宏用法/门谱/坑）。
% 键账：件内 ansblock 键集 ≡ 题面库片键集 ≡ manifest 键序（21 键，零缺漏零多余）；
%   印面号 1—16＝基础回归＋压轴收尾（bank 序直排），17—21＝课堂评价 G1—G5；键↔号映射见 值台账-课时19.json。
% 结构制式：\jietitle{章末总结与复习}＝总结提升（知识导学填空零键＋判断正误×2＋留白）→
%   基础回归（教材复习题A组，印1—6＝19-01~06，2题组）→ 压轴收尾（2.8压轴池，印7—16＝19-07~16，
%   3题组，每题组尾 \xiaojie＋2\bindp 方法小结）→ ketang 课堂评价（印17—21＝G1~G5）。
% 多选门：本片多选 2 道（19-11、19-12，印11/12）≤4 合规，成卷未增；印面带 \duoxuan\hspace{0.6em} 标记。
% 印面符号口径：源详解 ⇒ 箭头改「得/即」文读、✓/✗ 改「符合/不合」文读、〔亲算印证〕〔勘误留痕〕
%   〔源订正〕〔闸注〕源注不入印面；题面库【注】行为元数据不入题面；19-10/19-15 题面空位照源直排
%   \kongbai（19-10「如图」系冻结题面原文，本件无源图尺寸可回补，存照）；19-03 勘误依教材原文
%   （椭圆右半，定稿 §三.2）；19-G3 源订正 x−y+3=0（源 x−y+2=0 不过 P(1,4)，答案不受影响）；
%   19-13 Δ=240 照答案侧（亲算②-3 括注 Δ=60 系笔误，不采）。
% 值字节锁：\ansitem 值串自答案侧逐字回嵌零改动；19-16 值行含文本态下划线（x_A），经
%   \catcode 组内局部化照字印出（18-03 同法）。
% 括线判据（承重墙禁放宽）：估高＞8 行（wlen/23 栏宽口径）→ 括线模，逐键读数入值台账；
%   其余灰底。本片括线键以门谱读数为准（见值台账判模口径）。
% 门谱读数：见 过程件 工作区/_tmpM3S2波1补位4b_0914/（编译 log、门谱输出、台账）。
% ============================================================
\documentclass[fontset=none]{ctexart}
\usepackage{qp-m3}
% —— 印面逐字符号路由补钉（探针实证 0914：书宋缺 U+2212，▱ 诸字库皆缺→NSC 子块；
%     禁改 sty 本体保 md5；无 ▱/− 用件的片可整块照抄，零副作用） ——
\xeCJKDeclareCharClass{Default}{"2212}%
\setCJKfamilyfont{nscblk}[Path=C:/提示词/工作区/字替对照-0909/variantF/fonts/,BoldFont=NSC-w600.ttf]{NSC-w500.ttf}%
\xeCJKDeclareSubCJKBlock{nscblk}{"25B1}%
\newcommand{\pxparallelogram}{{\CJKfamily{nscblk}▱}}%
% —— 断行微调（纯题档/详解档三零门：CJK 胶收缩 0.01→0.05em；应急伸展撤行回落 sty 默认
%     1em 三零仍守（母版 2026-09-14 实测），补丁最少化；Underfull 治理读数见过程件） ——
\xeCJKsetup{CJKglue={\hskip 0.10em plus 0.08em minus 0.05em}}%
% —— 页脚件名（qp-headfoot 复用入口） ——
\renewcommand{\qpzhangming}{第二章\quad 平面解析几何}%
\renewcommand{\qpjianming}{导学件}%

\begin{document}

\zhangtitle{第二章\quad 平面解析几何}
\jietitle{章末总结与复习}
\mubiaoline
\mubiaomu{1}{系统梳理椭圆、双曲线、抛物线的定义、标准方程与几何性质，能准确辨析离心率范围、渐近线与准线等要素，会由几何条件反求参数或方程；}
\mubiaomu{2}{巩固直线与圆锥曲线位置关系的判定方法，规范处理联立消元、判别式讨论与退化情形（平行渐近线、平行对称轴），熟练运用弦长公式与点差法解决弦长与中点弦问题；}
\mubiaomu{3}{综合运用定义法、方程思想与设参消参策略处理定点、定值、最值与范围问题，体会数形结合与转化化归思想在压轴综合问题中的运用价值.}

\vspace{1.7mm}
\begin{multicols}{2}
\emergencystretch=1em

% ============================================================
% 总结提升（知识导学填空；零键零答案——\kongbai 空线＋\zhentib 判断题面形，
% 两档等形不泄答）
% ============================================================
\huaxing{总}{结}{提}{升}{知识结构\quad 通法提炼}

\zsd{一}{椭圆、双曲线、抛物线的定义与方程}

\tiaomu{1}{椭圆与双曲线的定义：平面内与两个定点\(F_1\)、\(F_2\)的距离之和为常数（大于\(|F_1F_2|\)）的点的轨迹是\kongbai{}；与两定点距离之差的绝对值为常数（小于\(|F_1F_2|\)且大于零）的点的轨迹是\kongbai{}．\zhuzhu{定义式中的常数范围是定义的组成部分，“大于/小于两定点间距离”不可省略；抛物线是平面内到定点与到定直线（准线）距离相等的点的轨迹，离心率恒为 1．}}

\tiaomu{2}{标准方程定位：椭圆\(x^{2}/a^{2}+y^{2}/b^{2}=1\)（\(a>b>0\)）的焦点在\kongbai{}轴上；双曲线\(x^{2}/a^{2}-y^{2}/b^{2}=1\)的渐近线方程为\kongbai{}；抛物线\(y^{2}=2px\)（\(p>0\)）的焦点坐标为\kongbai{}，准线方程为\kongbai{}．\zhuzhu{椭圆看分母大小定焦点位置，双曲线看平方项系数的符号定实轴位置；由离心率、渐近线、准线等条件求方程时，先定位再定量，防焦点位置误设．}}

\zsd{二}{直线与圆锥曲线的位置关系}

\tiaomu{1}{判定路径：联立直线与曲线方程消元，若得二次方程，则\(\Delta>0\)、\(\Delta=0\)、\(\Delta<0\)分别对应\kongbai{}；消元后为一次方程时（直线平行于双曲线的渐近线或抛物线的对称轴）恰有\kongbai{}个公共点．\zhuzhu{“相切”与“恰有一个公共点”不等价；二次项系数含参数时，先讨论系数为零的退化情形，再用判别式判型．}}

\tiaomu{2}{弦长与中点弦：弦长公式\(|AB|=\)\kongbai{}（\(k\) 存在）；中点弦用点差法：把两端点坐标代入曲线方程后相减，得斜率与中点坐标的关系，所求斜率须回验\kongbai{}．\zhuzhu{点差法解出的斜率必须代回联立方程检验判别式，防“无弦之弦”；抛物线的弦还要单列垂直于对称轴的情形．}}

\zsd{三}{通法综合：设参、化归与回验}

\tiaomu{1}{设线设参与消参：过定点的直线可设\(y=kx+m\)或\(x=my+n\)，垂直于坐标轴的情形单列；把垂直、等长、定比等几何条件译成坐标等式后，用根与系数的关系整体代换，使参数\kongbai{}即得定值、定点；最值与范围问题先由判别式与位置关系确定参数的\kongbai{}，再求目标函数的最值并检验取等．\zhuzhu{结论必须带回原约束核验：越界根舍、退化情形单列交代；开放条件题各路独立求解，殊途同归后共解后继问题．}}

\zhenhead{判断正误(正确的打\gou{},错误的打\cha{})}

\zhentib{(1)}{平面内到两定点距离之差的绝对值等于常数的点的轨迹一定是双曲线．}

\zhentib{(2)}{直线与双曲线的渐近线平行时，直线与该双曲线恰有一个公共点，且相切．}

\liubai[9mm]{此处书写}

% ============================================================
% 基础回归（教材复习题A组＝练槽 19-01~06；题后紧跟答案＝ansblock 内嵌，
% 值照答案侧逐字，详解承批E §七【详解】回嵌＋印面清洗）
% ============================================================
\huaxing{基}{础}{回}{归}{教材复习题A组\quad 夯基固本}

% ---------- 题组一 概念辨析与轨迹方程（印面 1—4） ----------
\tjdnr{一}{概念辨析与轨迹方程}{例\textbf{1}}{简单(知识点一)}{方程\(x^{2}-4x+1=0\)的两个根可分别作为\nobreak(\kongwei)}

\lxopt{A．椭圆和双曲线的离心率}

\lxopt{B．两抛物线的离心率}

\lxopt{C．椭圆和抛物线的离心率}

\lxopt{D．两椭圆的离心率}

\begin{ansblock}[2章-练-课时19-01]
% ans:2章-练-课时19-01
\ansitem{1}{A}
\ansnote{详解}{求根公式：\(x=2\pm\sqrt{3}\)，其中 \(2-\sqrt{3}\approx 0.27\in(0,1)\) 可作椭圆离心率，\(2+\sqrt{3}\approx 3.73\in(1,+\infty)\) 可作双曲线离心率．抛物线离心率恒等于 1，两根均不为 1，排除 B、C；两椭圆离心率都须在 \((0,1)\) 内，而 \(2+\sqrt{3}>1\)，排除 D．故选：A．}
\end{ansblock}

\liB{变式\textbf{2}}{简单(知识点一)}{}{曲线\(x^{2}/25+y^{2}/9=1\)与曲线\(x^{2}/(25-k)+y^{2}/(9-k)=1\)（\(k<9\)）的\nobreak(\kongwei)}

\lxopt{A．长轴长相等}

\lxopt{B．焦距相等}

\lxopt{C．离心率相等}

\lxopt{D．短轴长相等}

\begin{ansblock}[2章-练-课时19-02]
% ans:2章-练-课时19-02
\ansitem{2}{B}
\ansnote{详解}{\(k<9\) 得 \(25-k>9-k>0\)，第二曲线为椭圆，其 \(c^{2}=(25-k)-(9-k)=16\)，与 \(k\) 无关；第一曲线 \(c^{2}=25-9=16\)．故两曲线半焦距均为 4、焦距均为 8，选：B．长轴 10 与 \(2\sqrt{25-k}\)、短轴 6 与 \(2\sqrt{9-k}\)、离心率的平方 \(16/25\) 与 \(16/(25-k)\) 均随 \(k\) 变动，排除 A、C、D．}
\end{ansblock}

\tjdnr{一}{概念辨析与轨迹方程}{例\textbf{3}}{中档(知识点一)}{已知\(\triangle ABC\)的三边\(AB\)、\(BC\)、\(CA\)满足\(2|BC|=|AB|+|CA|\)，\(|AB|>|CA|\)，且\(B(-1,0)\)、\(C(1,0)\)，求点\(A\)的轨迹方程，并说明它是什么曲线．}
\xiexwei{14mm}

{\ansblockgrayfalse
\begin{ansblock}[2章-练-课时19-03]
% ans:2章-练-课时19-03
\ansitem{3}{x²/4+y²/3=1（x>0，y≠0），椭圆 x²/4+y²/3=1 的右半部分（不含与 x 轴的交点）}
\ansnote{详解}{\(|BC|=2\)，故 \(|AB|+|CA|=2|BC|=4>|BC|=2\)，由椭圆定义，\(A\) 的轨迹是以 \(B\)、\(C\) 为焦点、长半轴 \(a=2\)、半焦距 \(c=1\)、短半轴 \(b^{2}=a^{2}-c^{2}=3\) 的椭圆 \(x^{2}/4+y^{2}/3=1\)（去掉长轴端点，因 \(4>|BC|\) 严格成立）．在该椭圆上 \(|AB|=a+ex\)、\(|CA|=a-ex\)（\(e=1/2\)），\(|AB|>|CA|\) 即 \(x>0\)．又 \(A\) 为三角形顶点，不能与 \(B\)、\(C\) 共线，即 \(y\neq 0\)（舍去点 \((2,0)\)）．故轨迹为 \(x^{2}/4+y^{2}/3=1\)（\(x>0\)，\(y\neq 0\)），是椭圆的右半部分（不含与 \(x\) 轴的交点）．}
\end{ansblock}}

\liB{变式\textbf{4}}{中档(知识点一)}{}{若方程\((3-a)x^{2}+(a+1)y^{2}+a^{2}-2a-3=0\)是双曲线的方程，求实数\(a\)的取值范围．}
\xiexwei{14mm}

\begin{ansblock}[2章-练-课时19-04]
% ans:2章-练-课时19-04
\ansitem{4}{a<−1 或 a>3}
\ansnote{详解}{双曲线型须二次项系数异号：\((3-a)(a+1)<0\)，即 \(a<-1\) 或 \(a>3\)．此时常数项 \(a^{2}-2a-3=(a-3)(a+1)>0\) 自动相容，方程确为双曲线：\(a>3\) 时移项配方得 \(x^{2}/(a+1)-y^{2}/(a-3)=1\)；\(a<-1\) 时同法得 \(x^{2}/(a+1)-y^{2}/(a-3)=1\)（分子分母同乘 \(-1\) 后仍为平方差形式）．即 \(-1<a<3\) 时为椭圆型或虚轨迹、退化，不取．故 \(a\in(-\infty,-1)\cup(3,+\infty)\)．}
\end{ansblock}

% ---------- 题组二 位置关系与公共点（印面 5—6） ----------
\tjdnr{二}{位置关系与公共点}{例\textbf{5}}{中档(知识点二)}{已知直线\(y=kx+2\)与椭圆\(x^{2}+2y^{2}=2\)相交于不同的两点，求\(k\)的取值范围．}
\xiexwei{14mm}

\begin{ansblock}[2章-练-课时19-05]
% ans:2章-练-课时19-05
\ansitem{5}{k<−√6/2 或 k>√6/2}
\ansnote{详解}{联立消 \(y\)：\((1+2k^{2})x^{2}+8kx+6=0\)．相交于不同的两点即 \(\Delta>0\)：\(\Delta=64k^{2}-24(1+2k^{2})=16k^{2}-24=8(2k^{2}-3)>0\)，即 \(k^{2}>3/2\)，\(|k|>\sqrt{6}/2\)．故 \(k<-\sqrt{6}/2\) 或 \(k>\sqrt{6}/2\)．}
\end{ansblock}

\liB{变式\textbf{6}}{简单(知识点二)}{}{直线\(l\)过点\(P(2,4)\)且与抛物线\(y^{2}=8x\)只有一个公共点，求直线\(l\)的方程．}
\xiexwei{14mm}

{\ansblockgrayfalse
\begin{ansblock}[2章-练-课时19-06]
% ans:2章-练-课时19-06
\ansitem{6}{y=4 或 y=x+2}
\ansnote{详解}{①斜率不存在：\(x=2\) 代入 \(y^{2}=16\) 得两点 \((2,\pm 4)\)，不合，舍．②斜率 \(k=0\)：\(y=4\)，与抛物线恰一公共点 \((2,4)\)，符合．③\(k\neq 0\)：设 \(y-4=k(x-2)\)，以 \(x=y^{2}/8\) 代入并整理：\(ky^{2}-8y+32-16k=0\)，恰一公共点即 \(\Delta=64-4k(32-16k)=64(k-1)^{2}=0\)，得 \(k=1\)，即 \(y=x+2\)．综上共两条：\(y=4\) 与 \(y=x+2\)．易错提醒：与对称轴平行的情形（\(k=0\)）不可漏．}
\end{ansblock}}

% ============================================================
% 压轴收尾（2.8压轴池＝练槽 19-07~16；3题组，每题组尾 \xiaojie＋2\bindp 方法小结）
% ============================================================
\huaxing{压}{轴}{收}{尾}{2.8压轴池\quad 能力跨越}

% ---------- 题组一 方程直求与离心率（印面 7—10） ----------
\tjdnr{一}{方程直求与离心率}{例\textbf{7}}{中档(知识点一)}{双曲线\(C\)：\(x^{2}/a^{2}-y^{2}/b^{2}=1\)（\(a>0\)，\(b>0\)）的离心率为\(\sqrt{2}\)，抛物线\(y^{2}=2px\)（\(p>0\)）的准线与双曲线\(C\)的渐近线交于\(A\)、\(B\)两点，\(\triangle OAB\)（\(O\)为坐标原点）的面积为4，则抛物线的方程为\nobreak(\kongwei)}

\bindopt {\fontsize{10.5pt}{\qpTJgLeadLiti}\selectfont \makebox[\dimexpr0.5\linewidth-1em\relax][l]{A．\(y^{2}=4x\)}\makebox[\dimexpr0.5\linewidth-1em\relax][l]{B．\(y^{2}=6x\)}\par}

\bindopt {\fontsize{10.5pt}{\qpTJgLeadLiti}\selectfont \makebox[\dimexpr0.5\linewidth-1em\relax][l]{C．\(y^{2}=8x\)}\makebox[\dimexpr0.5\linewidth-1em\relax][l]{D．\(y^{2}=16x\)}\par}

\begin{ansblock}[2章-练-课时19-07]
% ans:2章-练-课时19-07
\ansitem{7}{C}
\ansnote{详解}{\(e=\sqrt{2}\) 得 \(a=b\)，双曲线为等轴双曲线，渐近线 \(y=\pm x\)．准线 \(x=-p/2\)，交渐近线于 \(A(-p/2,p/2)\)、\(B(-p/2,-p/2)\)，\(|AB|=p\)，\(O\) 到 \(AB\) 的距离为 \(p/2\)，故 \(S=\frac{1}{2}\cdot p\cdot\frac{p}{2}=p^{2}/4=4\)，得 \(p=4\)，抛物线的方程为 \(y^{2}=8x\)．故选：C．}
\end{ansblock}

\liB{变式\textbf{8}}{难(知识点二)}{}{已知双曲线\(\Gamma\)：\(x^{2}-y^{2}=4\)，双曲线\(\Gamma\)的右焦点为\(F\)，圆\(C\)的圆心在\(y\)轴正半轴上，且经过坐标原点\(O\)，圆\(C\)与双曲线\(\Gamma\)的右支交于\(A\)、\(B\)两点．}

\bindp (1) 当\(\triangle OFA\)是以\(F\)为直角顶点的直角三角形时，求\(\triangle OFA\)的面积；

\bindp (2) 若点\(A\)的坐标是\((\sqrt{5},1)\)，求直线\(AB\)的方程；

\bindp (3) 求证：直线\(AB\)与圆\(x^{2}+y^{2}=2\)相切．
\xiexwei{18mm}

{\ansblockgrayfalse
\begin{ansblock}[2章-练-课时19-08]
% ans:2章-练-课时19-08
\ansitem{8}{(1) 2√2；(2) y=[(2√2+√5)/3]x−(2√10+2)/3；(3) 证毕（O 到 AB 距离恒为 √2）}
\ansnote{详解}{\(c^{2}=4+4=8\)，故 \(F(2\sqrt{2},0)\)．圆心 \((0,m)\)（\(m>0\)）且过 \(O\)，半径为 \(m\)，圆 \(C\)：\(x^{2}+y^{2}-2my=0\)．\par\noindent (1) \(AF\perp OF\)（\(x\) 轴）得 \(x_A=2\sqrt{2}\)，代入 \(\Gamma\) 得 \(y_A^{2}=8-4=4\)，取 \(A(2\sqrt{2},2)\)，故 \(S=\frac{1}{2}\cdot|OF|\cdot|y_A|=\frac{1}{2}\cdot 2\sqrt{2}\cdot 2=2\sqrt{2}\)．\par\noindent (2) 代 \(A(\sqrt{5},1)\)：\(5+1-2m=0\)，得 \(m=3\)．圆与 \(\Gamma\) 联立（\(x^{2}-y^{2}=4\) 与 \(x^{2}+y^{2}-6y=0\) 相减）得 \(y^{2}-3y+2=0\)，即 \(y=1\)（点 \(A\)）或 \(y=2\)，故 \(B(2\sqrt{2},2)\)．\(k_{AB}=(2-1)/(2\sqrt{2}-\sqrt{5})=(2\sqrt{2}+\sqrt{5})/3\)（分母有理化，\((2\sqrt{2})^{2}-(\sqrt{5})^{2}=3\)）；截距 \(=1-k\sqrt{5}=1-(2\sqrt{10}+5)/3=-(2\sqrt{10}+2)/3\)，故直线 \(AB\)：\(y=\frac{2\sqrt{2}+\sqrt{5}}{3}x-\frac{2\sqrt{10}+2}{3}\)．\par\noindent (3) 一般情形：圆 \(x^{2}+y^{2}-2my=0\) 与 \(\Gamma\) 消 \(x\) 得 \(y^{2}-my+2=0\)，故 \(y_1y_2=2\)、\(y_1^{2}+y_2^{2}=m^{2}-4\)；右支上 \(x_i=\sqrt{y_i^{2}+4}\)，故 \(x_1x_2=\sqrt{(y_1y_2)^{2}+4(y_1^{2}+y_2^{2})+16}=2\sqrt{m^{2}+1}\)、\(x_1^{2}+x_2^{2}=m^{2}+4\)．原点到 \(AB\) 的距离 \(d=|x_1y_2-x_2y_1|/|AB|\)，其中分子平方 \(=x_1^{2}y_2^{2}+x_2^{2}y_1^{2}-2x_1x_2y_1y_2=y_2^{2}(y_1^{2}+4)+y_1^{2}(y_2^{2}+4)-4\sqrt{m^{2}+1}\cdot 2=2(y_1y_2)^{2}+4(y_1^{2}+y_2^{2})-8\sqrt{m^{2}+1}=4m^{2}-8-8\sqrt{m^{2}+1}\)；分母平方 \(=(x_1^{2}+x_2^{2}-2x_1x_2)+(y_1^{2}+y_2^{2}-2y_1y_2)=(m^{2}+4-4\sqrt{m^{2}+1})+(m^{2}-4-4)=2m^{2}-4-4\sqrt{m^{2}+1}\)．分子平方 \(=2\times\)分母平方恒成立（\(m>2\sqrt{2}\) 时分母 \(>0\)），故 \(d^{2}\equiv 2\)，即 \(d=\sqrt{2}\) 恒等于圆 \(x^{2}+y^{2}=2\) 的半径，直线 \(AB\) 与该圆相切，证毕．}
\end{ansblock}}

\tjdnr{一}{方程直求与离心率}{例\textbf{9}}{中档(知识点一)}{设\(F_1\)、\(F_2\)是椭圆\(C_1\)：\(x^{2}/a_1^{2}+y^{2}/b_1^{2}=1\)（\(a_1>b_1>0\)）与双曲线\(C_2\)：\(x^{2}/a_2^{2}-y^{2}/b_2^{2}=1\)（\(a_2>0\)，\(b_2>0\)）的公共焦点，曲线\(C_1\)、\(C_2\)在第一象限内交于点\(M\)，\(\angle F_1MF_2=90^{\circ}\)．若椭圆的离心率\(e_1\in[\sqrt{6}/3,1)\)，则双曲线的离心率\(e_2\)的范围是\nobreak(\kongwei)}

\bindopt {\fontsize{10.5pt}{\qpTJgLeadLiti}\selectfont \makebox[\dimexpr0.5\linewidth-1em\relax][l]{A．\((1,\sqrt{2}]\)}\makebox[\dimexpr0.5\linewidth-1em\relax][l]{B．\((1,\sqrt{3}]\)}\par}

\bindopt {\fontsize{10.5pt}{\qpTJgLeadLiti}\selectfont \makebox[\dimexpr0.5\linewidth-1em\relax][l]{C．\([\sqrt{3},+\infty)\)}\makebox[\dimexpr0.5\linewidth-1em\relax][l]{D．\([\sqrt{2},+\infty)\)}\par}

{\ansblockgrayfalse
\begin{ansblock}[2章-练-课时19-09]
% ans:2章-练-课时19-09
\ansitem{9}{A}
\ansnote{详解}{设 \(|MF_1|>|MF_2|\)，则 \(|MF_1|+|MF_2|=2a_1\)、\(|MF_1|-|MF_2|=2a_2\)，得 \(|MF_1|=a_1+a_2\)、\(|MF_2|=a_1-a_2\)．\(\angle F_1MF_2=90^{\circ}\) 得 \(|MF_1|^{2}+|MF_2|^{2}=4c^{2}\)，即 \((a_1+a_2)^{2}+(a_1-a_2)^{2}=4c^{2}\)，故 \(a_1^{2}+a_2^{2}=2c^{2}\)，同除 \(c^{2}\)：\(1/e_1^{2}+1/e_2^{2}=2\)．\(e_1\in[\sqrt{6}/3,1)\) 得 \(1/e_1^{2}\in(1,3/2]\)，故 \(1/e_2^{2}=2-1/e_1^{2}\in[1/2,1)\)，即 \(e_2^{2}\in(1,2]\)，\(e_2\in(1,\sqrt{2}]\)．故选：A．}
\end{ansblock}}

\liB{变式\textbf{10}}{中档(知识点一)}{}{如图，\(F_1\)、\(F_2\)是椭圆\(C_1\)与双曲线\(C_2\)的公共焦点，\(A\)、\(B\)分别是\(C_1\)、\(C_2\)在第二、四象限的公共点，若\(|OF_1|=\frac{1}{2}|AB|\)，且\(\angle OF_1B=\frac{\pi}{6}\)，则\(C_1\)与\(C_2\)的离心率之积为\kongbai{}．}
\xiexwei{10mm}

{\ansblockgrayfalse
\begin{ansblock}[2章-练-课时19-10]
% ans:2章-练-课时19-10
\ansitem{10}{2}
\ansnote{详解}{两曲线关于原点对称，故 \(O\) 为 \(AB\) 中点，四边形 \(AF_1BF_2\) 为平行四边形，其对角线 \(|AB|=2|OF_1|=2c\) 与 \(|F_1F_2|=2c\) 相等，故四边形为矩形，\(\angle F_1AF_2=\angle F_1BF_2=90^{\circ}\)．\(\triangle OF_1B\) 中 \(|OF_1|=|OB|=c\)，\(\angle OF_1B=\pi/6\)，顶角 \(\angle F_1OB=2\pi/3\)；由 \(A\) 与 \(B\) 关于原点对称得 \(\angle F_1OA=\pi-2\pi/3=\pi/3\)，而 \(|OA|=|OF_1|=c\)，故 \(\triangle AOF_1\) 为等边三角形，\(|AF_1|=c\) 且 \(\angle AF_1O=\pi/3\)，\(\angle AF_1B=\pi/3+\pi/6=\pi/2\)，与矩形一致．于是 \(|BF_1|=\sqrt{|AB|^{2}-|AF_1|^{2}}=\sqrt{4c^{2}-c^{2}}=\sqrt{3}c\)，且平行四边形对边 \(|AF_2|=|BF_1|=\sqrt{3}c\)．椭圆定义：\(c+\sqrt{3}c=2a_1\)，故 \(e_1=2/(\sqrt{3}+1)\)；双曲线定义：\(\sqrt{3}c-c=2a_2\)，故 \(e_2=2/(\sqrt{3}-1)\)．\(e_1e_2=4/[(\sqrt{3}+1)(\sqrt{3}-1)]=4/2=2\)．}
\end{ansblock}}

\xiaojie{①识别：以准线截渐近线的面积定方程、等轴双曲线、共焦点张角直角的 \(1/e_1^{2}+1/e_2^{2}=2\) 型恒等式、焦点四边形的矩形化等为目标的题，判归本题型．}
\bindp {\kaishu ②操作：离心率先定型（椭圆 \(<1\)、双曲线 \(>1\)、抛物线 \(=1\)）；共焦点直角用定义联立焦半径出恒等式；几何条件（平行四边形、矩形、等边三角形）先把角度与边长定死，再代定义求基本量．}
\bindp {\kaishu ③收束：方程直求后回验定位（焦点位置、开口方向、准线归属）；范围结论核对端点开闭与取等情形；三问压轴第(3)问常由特例过渡到一般恒等式．}

% ---------- 题组二 曲线族多选与中点弦（印面 11—13） ----------
\tjdnr{二}{曲线族多选与中点弦}{例\textbf{11}}{中档(知识点一)}{\duoxuan\hspace{0.6em}曲线\(C\)的方程为\(x^{2}/(16+\lambda)+y^{2}/(9+\lambda)=1\)，则下列说法正确的是\nobreak(\kongwei)}

\lxopt{A．存在实数\(\lambda\)使得曲线\(C\)的轨迹为圆}

\lxopt{B．存在实数\(\lambda\)使得曲线\(C\)的轨迹为椭圆}

\lxopt{C．存在实数\(\lambda\)使得曲线\(C\)的轨迹为双曲线}

\lxopt{D．无论\(\lambda\)（\(\lambda>-16\)且\(\lambda\neq-9\)）取何值，曲线\(C\)的焦距为定值}

\begin{ansblock}[2章-练-课时19-11]
% ans:2章-练-课时19-11
\ansitem{11}{BCD}
\ansnote{详解}{A：圆须 \(16+\lambda=9+\lambda\)，矛盾，不存在，不合．B：\(\lambda>-9\) 时 \(16+\lambda>9+\lambda>0\)，方程表示焦点在 \(x\) 轴上的椭圆，符合．C：\(-16<\lambda<-9\) 时 \((16+\lambda)(9+\lambda)<0\)，表示焦点在 \(x\) 轴上的双曲线，符合．D：\(\lambda>-9\) 时为椭圆，\(c^{2}=(16+\lambda)-(9+\lambda)=7\)；\(-16<\lambda<-9\) 时为双曲线，\(c^{2}=(16+\lambda)+[-(9+\lambda)]=7\)，焦距恒为 \(2\sqrt{7}\)，符合．故选：BCD．}
\end{ansblock}

\liB{变式\textbf{12}}{中档(知识点一)}{}{\duoxuan\hspace{0.6em}已知椭圆\(C_1\)：\(x^{2}/a^{2}+y^{2}/b^{2}=1\)（\(a>b>0\)）的左、右焦点分别为\(F_1\)、\(F_2\)，离心率为\(e_1\)，椭圆\(C_1\)的上顶点为\(M\)，且\(\overrightarrow{MF_1}\cdot\overrightarrow{MF_2}=0\)，双曲线\(C_2\)和椭圆\(C_1\)有相同的焦点，且双曲线\(C_2\)的离心率为\(e_2\)，\(P\)为曲线\(C_1\)与\(C_2\)的一个公共点．若\(\angle F_1PF_2=\pi/3\)，则\nobreak(\kongwei)}

\lxopt{A．\(e_2/e_1=2\)}

\lxopt{B．\(e_1e_2=\sqrt{3}/2\)}

\lxopt{C．\(e_1^{2}+e_2^{2}=5/2\)}

\lxopt{D．\(e_2^{2}-e_1^{2}=1\)}

{\ansblockgrayfalse
\begin{ansblock}[2章-练-课时19-12]
% ans:2章-练-课时19-12
\ansitem{12}{BD}
\ansnote{详解}{\(\overrightarrow{MF_1}\cdot\overrightarrow{MF_2}=0\) 且 \(|MF_1|=|MF_2|=\sqrt{b^{2}+c^{2}}\)，故 \(\triangle MF_1F_2\) 为等腰直角三角形，\(b=c\)，故 \(e_1^{2}=b^{2}/(b^{2}+c^{2})=1/2\)，\(e_1=\sqrt{2}/2\)，且 \(a=\sqrt{2}c\)．在焦点三角形 \(\triangle PF_1F_2\) 中 \(\angle F_1PF_2=\pi/3\)，设 \(|PF_1|=m\)、\(|PF_2|=n\)（\(m>n\)，\(P\) 取双曲线右支），则 \(m+n=2a_1=2\sqrt{2}c\)、\(m-n=2a_2\)．余弦定理：\(4c^{2}=m^{2}+n^{2}-2mn\cos(\pi/3)=(m+n)^{2}-3mn\)，故 \(mn=(8c^{2}-4c^{2})/3=4c^{2}/3\)．于是 \((m-n)^{2}=(m+n)^{2}-4mn=8c^{2}-16c^{2}/3=8c^{2}/3\)，即 \(4a_2^{2}=8c^{2}/3\)，\(a_2^{2}=2c^{2}/3\)，\(e_2^{2}=3/2\)，\(e_2=\sqrt{6}/2\)．逐项：\(e_1e_2=(\sqrt{2}/2)\cdot(\sqrt{6}/2)=\sqrt{3}/2\)，B 符合；\(e_2^{2}-e_1^{2}=3/2-1/2=1\)，D 符合；\(e_2/e_1=\sqrt{3}\neq 2\)，A 不合；\(e_1^{2}+e_2^{2}=2\neq 5/2\)，C 不合．故选：BD．}
\end{ansblock}}

\tjdnr{二}{曲线族多选与中点弦}{例\textbf{13}}{中档(知识点二)}{已知双曲线\(C\)：\(x^{2}/a^{2}-y^{2}/b^{2}=1\)（\(a>0\)，\(b>0\)）与椭圆\(x^{2}/18+y^{2}/14=1\)有共同的焦点，点\(A(3,\sqrt{7})\)在双曲线\(C\)上．}

\bindp (1) 求双曲线\(C\)的方程及渐近线方程；

\bindp (2) 以\(P(1,2)\)为中点作双曲线\(C\)的一条弦\(AB\)，求弦\(AB\)所在直线的方程．
\xiexwei{18mm}

{\ansblockgrayfalse
\begin{ansblock}[2章-练-课时19-13]
% ans:2章-练-课时19-13
\ansitem{13}{(1) x²/2−y²/2=1、y=±x；(2) x−2y+3=0}
\ansnote{详解}{(1) 椭圆 \(c^{2}=18-14=4\)，故双曲线 \(c=2\)、\(a^{2}+b^{2}=4\)；\(A\) 代入：\(9/a^{2}-7/b^{2}=1\)．令 \(u=a^{2}\)：\(9/u-7/(4-u)=1\)，即 \(9(4-u)-7u=u(4-u)\)，整理得 \(u^{2}-20u+36=0\)，解得 \(u=2\) 或 \(u=18\)（\(>c^{2}\) 舍），故 \(a^{2}=b^{2}=2\)，双曲线 \(C\)：\(x^{2}/2-y^{2}/2=1\)，渐近线 \(y=\pm x\)（等轴双曲线）．\par\noindent (2) 点差法：\(x_1^{2}-y_1^{2}=2\)、\(x_2^{2}-y_2^{2}=2\) 相减得 \((x_1+x_2)(x_1-x_2)=(y_1+y_2)(y_1-y_2)\)，故 \(k_{AB}=(y_1-y_2)/(x_1-x_2)=(x_1+x_2)/(y_1+y_2)=2/4=1/2\)，直线为 \(y-2=\frac{1}{2}(x-1)\)，即 \(x-2y+3=0\)．回验：联立消 \(y\) 得 \(3x^{2}-6x-17=0\)，\(\Delta=36+204=240>0\)，为真弦，符合．}
\end{ansblock}}

\xiaojie{①识别：含参曲线族逐项判型、共焦点张角（直角、\(60^{\circ}\)）联立定义与余弦定理、中点弦点差法等目标的题，判归本题型．}
\bindp {\kaishu ②操作：曲线族按参数分段讨论分母符号定型；焦点三角形用定义式＋余弦定理联出基本量；点差法斜率求出后回验判别式，防“无弦之弦”．}
\bindp {\kaishu ③收束：多选逐项独立验真伪，反例与特值优先；中点弦回验交点为不同两实点，退化与越界根（\(a^{2}>c^{2}\)）单列舍取．}

% ---------- 题组三 文化载体与压轴综合（印面 14—16） ----------
\tjdnr{三}{文化载体与压轴综合}{例\textbf{14}}{中档(知识点一)}{黄金分割起源于公元前6世纪古希腊的毕达哥拉斯学派，公元前4世纪，古希腊数学家欧多克索斯第一个系统研究了这一问题，公元前300年前后欧几里得撰写《几何原本》时吸收了欧多克索斯的研究成果，进一步系统论述了黄金分割，成为最早的有关黄金分割的论著。黄金分割是指将整体一分为二，较大部分与整体部分的比值等于较小部分与较大部分的比值，其比值为\((\sqrt{5}-1)/2\)，把\((\sqrt{5}-1)/2\)称为黄金分割数．已知双曲线\(x^{2}/(\sqrt{5}-1)^{2}-y^{2}/m=1\)的实轴长与焦距的比值恰好是黄金分割数，则\(m\)的值为\nobreak(\kongwei)}

\bindopt {\fontsize{10.5pt}{\qpTJgLeadLiti}\selectfont \makebox[\dimexpr0.25\linewidth-0.5em\relax][l]{A．\(2\sqrt{5}-2\)}\makebox[\dimexpr0.25\linewidth-0.5em\relax][l]{B．\(\sqrt{5}+1\)}\makebox[\dimexpr0.25\linewidth-0.5em\relax][l]{C．\(2\)}\makebox[\dimexpr0.25\linewidth-0.5em\relax][l]{D．\(2\sqrt{5}\)}\par}

\begin{ansblock}[2章-练-课时19-14]
% ans:2章-练-课时19-14
\ansitem{14}{A}
\ansnote{详解}{\(a^{2}=(\sqrt{5}-1)^{2}=6-2\sqrt{5}\)．\(a/c=(\sqrt{5}-1)/2\) 得 \(c=2a/(\sqrt{5}-1)=a(\sqrt{5}+1)/2\)，故 \(c^{2}=a^{2}(3+\sqrt{5})/2\)．\(m=c^{2}-a^{2}=a^{2}[(3+\sqrt{5})/2-1]=a^{2}(1+\sqrt{5})/2=(6-2\sqrt{5})(\sqrt{5}+1)/2=(3-\sqrt{5})(\sqrt{5}+1)=3\sqrt{5}+3-5-\sqrt{5}=2\sqrt{5}-2\)．故选：A．}
\end{ansblock}

\liB{变式\textbf{15}}{中档(知识点一)}{}{在①左顶点为\((-3,0)\)，②双曲线过点\((3\sqrt{2},4)\)，③离心率\(e=5/3\)这三个条件中任选一个，补充在下面问题中并作答（注：如果选择多个条件分别解答，按第一个解答计分）．}

\bindp 问题：已知双曲线与椭圆\(x^{2}/49+y^{2}/24=1\)共焦点，且\kongbai{}．

\bindp (1) 求双曲线的方程；

\bindp (2) 若点\(P\)在双曲线上，且\(|PF_1|=8\)，求\(|PF_2|\)．
\xiexwei{18mm}

{\ansblockgrayfalse
\begin{ansblock}[2章-练-课时19-15]
% ans:2章-练-课时19-15
\ansitem{15}{(1) 三条件均得 x²/9−y²/16=1；(2) |PF₂|=2 或 14}
\ansnote{详解}{(1) \(c^{2}=49-24=25\)，故 \(c=5\)．选①：左顶点 \((-3,0)\) 得 \(a=3\)．选②：\(18/a^{2}-16/b^{2}=1\) 配 \(a^{2}+b^{2}=25\)，即 \(18(25-a^{2})-16a^{2}=a^{2}(25-a^{2})\)，整理得 \(a^{4}-59a^{2}+450=0\)，解得 \(a^{2}=9\) 或 \(a^{2}=50\)（\(>25\) 舍），故 \(a^{2}=9\)．选③：\(e=c/a=5/3\)，即 \(5/a=5/3\)，得 \(a=3\)．三路均得 \(b^{2}=25-9=16\)，双曲线的方程为 \(x^{2}/9-y^{2}/16=1\)．\par\noindent (2) \(\bigl||PF_1|-|PF_2|\bigr|=2a=6\)，故 \(|PF_2|=2\) 或 \(14\)，均须回验：\(|PF_2|=2\) 时 \(P\) 为右顶点 \((3,0)\)（\(|PF_1|=3+5=8\)，且 \(8-2=6\) 合定义——\(P\)、\(F\) 成线段，退化但合法）；\(|PF_2|=14\) 时 \(P\) 在左支（\(x=-33/5\)，\(|PF_1|=8\)、\(14-8=6\)，且 \(8+14=22>2c=20\) 三角形存在）．故 \(|PF_2|=2\) 或 \(14\)．}
\end{ansblock}}

\tjdnr{三}{文化载体与压轴综合}{例\textbf{16}}{难(知识点二)}{如图所示，曲线\(C_1\)：\(y^{2}=2p_1x\)（\(p_1>0\)），曲线\(C_2\)：\(y^{2}=-2px\)（\(p>0\)），过点\(C(-1,0)\)作直线交曲线\(C_1\)于点\(A\)，交曲线\(C_2\)于点\(B\)，若点\(C\)在曲线\(C_1\)的准线上．}

\bindp (1) 求\(p_1\)；

\bindp (2) 若存在直线使点\(B\)为\(AC\)中点，求\(A\)点横坐标（用\(p\)表示）及\(AC\)斜率的范围．
\xiexwei{18mm}

{\ansblockgrayfalse
\begin{ansblock}[2章-练-课时19-16]
% ans:2章-练-课时19-16
{\catcode`\_=12 \ansitem{16}{(1) p₁=2；(2) x_A=p/(p+1)，斜率范围 (−1,0)∪(0,1)}}
\ansnote{详解}{(1) \(C_1\) 的准线 \(x=-p_1/2=-1\)，得 \(p_1=2\)，即 \(C_1\)：\(y^{2}=4x\)．\par\noindent (2) 设 \(B(-t^{2}/(2p),t)\)（\(C_2\) 上的参数点），\(B\) 为 \(AC\) 中点、\(C(-1,0)\)，故 \(A=2B-C=(1-t^{2}/p,2t)\)．\(A\) 在 \(C_1\) 上：\(4t^{2}=4(1-t^{2}/p)\)，得 \(t^{2}=p/(p+1)\in(0,1)\)（\(p>0\)），故 \(A\) 的横坐标 \(x_A=1-t^{2}/p=1-1/(p+1)=p/(p+1)\)（\(0<x_A<1\)，符合）．斜率：\(k_{AC}=k_{BC}=t/[-t^{2}/(2p)+1]\)；由 \(t^{2}=p/(p+1)\) 得 \(1/p=1/t^{2}-1\)，故分母 \(=-(t^{2}/2)(1/t^{2}-1)+1=(1+t^{2})/2\)，即 \(k=2t/(1+t^{2})\)（\(t\neq 0\)，否则 \(A=C\) 不成弦）．\(k^{2}=4t^{2}/(1+t^{2})^{2}\)，由 \(0<t^{2}<1\) 得 \(0<k^{2}<1\)（等号 \(t^{2}=1\) 不可达），故 \(k\in(-1,0)\cup(0,1)\)．}
\end{ansblock}}

\xiaojie{①识别：文化载体（比值定离心率）反求参数、三选一条件开放求方程、双抛物线中点跨支综合等形态，判归本题型．}
\bindp {\kaishu ②操作：文化题先把数量关系（比值、离心率）译成 \(a\)、\(c\) 的等式再代标准方程；开放题三路各设基本量独立求解，殊途同归后共解后问；参数点设元（\(B(-t^{2}/(2p),t)\)）把中点关系化为一元代入．}
\bindp {\kaishu ③收束：定义双解（\(|PF_2|=2\) 或 14）逐一回验定义式与三角形存在性，退化情形单列注记；斜率范围以参数平方设元求值域，端点与退化（\(t=0\)）单独交代．}

% ============================================================
% 课堂评价 G5（恒5＝3单选+2填空＝导槽 G1~G5；ketang 新制顺排可跨栏，禁 \vbox 装栏）
% 印面号 17—21；多选门：本片课堂评价无多选，全片 2（印11/12）≤4 合规
% ============================================================
\begin{ketang}{本组共5题（第17—21题），17—19为单选，20—21为填空．}

\jiancestem{{\fontsize{11.4pt}{14pt}\selectfont\heihao {\numboldjian 17}．}\hspace{4.1mm}\tieside{简单(知识点二)}\hspace{2.2mm}若直线\(l\)：\(y=kx+2\)与双曲线\(C\)：\(x^{2}-y^{2}=4\)的左、右两支各有一个交点，则实数\(k\)的取值范围是\nobreak(\kongwei)}

\bindopt {\fontsize{10.5pt}{\qpTJgLeadPingjia}\selectfont \makebox[\dimexpr0.5\linewidth-1em\relax][l]{A．\((-2,-1)\)}\makebox[\dimexpr0.5\linewidth-1em\relax][l]{B．\((1,2)\)}\par}

\bindopt {\fontsize{10.5pt}{\qpTJgLeadPingjia}\selectfont \makebox[\dimexpr0.5\linewidth-1em\relax][l]{C．\((-2,2)\)}\makebox[\dimexpr0.5\linewidth-1em\relax][l]{D．\((-1,1)\)}\par}

{\ansblockgrayfalse
\begin{ansblock}[2章-导-课时19-G1]
% ans:2章-导-课时19-G1
\ansitem{17}{D}
\ansnote{详解}{\(C\) 即 \(x^{2}/4-y^{2}/4=1\)，实轴在 \(x\) 轴上，渐近线 \(y=\pm x\)．联立 \(l\) 与 \(C\)：\(x^{2}-(kx+2)^{2}=4\)，即 \((1-k^{2})x^{2}-4kx-8=0\)．左支 \(x<0\)、右支 \(x>0\)，故「两支各一交点」等价于方程两实根异号：\(x_1x_2=-8/(1-k^{2})<0\)，即 \(1-k^{2}>0\)，得 \(-1<k<1\)；此时 \(\Delta=16k^{2}+32(1-k^{2})=32-16k^{2}>0\) 自动满足（\(k=\pm 1\) 时方程退化为一元一次、至多一交点，已排除）．几何印证：\(l\) 过定点 \((0,2)\)，\(k=\pm 1\) 时平行于渐近线仅扫一支．故 \(-1<k<1\)，选：D．}
\end{ansblock}}

\jiancestem{{\fontsize{11.4pt}{14pt}\selectfont\heihao {\numboldjian 18}．}\hspace{4.1mm}\tieside{简单(知识点一)}\hspace{2.2mm}若双曲线\(x^{2}/a^{2}-y^{2}/b^{2}=1\)（\(a>0\)，\(b>0\)）与直线\(y=\sqrt{3}x\)有交点，则其离心率的取值范围是\nobreak(\kongwei)}

\bindopt {\fontsize{10.5pt}{\qpTJgLeadPingjia}\selectfont \makebox[\dimexpr0.25\linewidth-0.5em\relax][l]{A．\((1,2)\)}\makebox[\dimexpr0.25\linewidth-0.5em\relax][l]{B．\((1,2]\)}\makebox[\dimexpr0.25\linewidth-0.5em\relax][l]{C．\((2,+\infty)\)}\makebox[\dimexpr0.25\linewidth-0.5em\relax][l]{D．\([2,+\infty)\)}\par}

\begin{ansblock}[2章-导-课时19-G2]
% ans:2章-导-课时19-G2
\ansitem{18}{C}
\ansnote{详解}{渐近线 \(y=\pm(b/a)x\)．直线 \(y=\sqrt{3}x\) 过原点：当 \(\sqrt{3}<b/a\) 时直线落在含 \(x\) 轴的开口区域内，与右支必有交点；当 \(\sqrt{3}=b/a\) 时直线即渐近线，无交点；当 \(\sqrt{3}>b/a\) 时直线落在含 \(y\) 轴的区域内，双曲线在该区域无点．故「有交点」等价于 \(b/a>\sqrt{3}\)．于是 \(e=\sqrt{1+b^{2}/a^{2}}>\sqrt{1+3}=2\)，即 \(e\in(2,+\infty)\)．故选：C．}
\end{ansblock}

\jiancestem{{\fontsize{11.4pt}{14pt}\selectfont\heihao {\numboldjian 19}．}\hspace{4.1mm}\tieside{简单(知识点二)}\hspace{2.2mm}已知直线\(l\)：\(x-y+3=0\)与双曲线\(C\)：\(x^{2}/a^{2}-y^{2}/b^{2}=1\)（\(a>0\)，\(b>0\)）交于\(A\)、\(B\)两点，点\(P(1,4)\)是弦\(AB\)的中点，则双曲线\(C\)的离心率为\nobreak(\kongwei)}

\bindopt {\fontsize{10.5pt}{\qpTJgLeadPingjia}\selectfont \makebox[\dimexpr0.25\linewidth-0.5em\relax][l]{A．\(4/3\)}\makebox[\dimexpr0.25\linewidth-0.5em\relax][l]{B．\(2\)}\makebox[\dimexpr0.25\linewidth-0.5em\relax][l]{C．\(\sqrt{5}/2\)}\makebox[\dimexpr0.25\linewidth-0.5em\relax][l]{D．\(\sqrt{5}\)}\par}

{\ansblockgrayfalse
\begin{ansblock}[2章-导-课时19-G3]
% ans:2章-导-课时19-G3
\ansitem{19}{D}
\ansnote{详解}{（点差法）设 \(A(x_1,y_1)\)、\(B(x_2,y_2)\)，由中点 \(P(1,4)\)：\(x_1+x_2=2\)、\(y_1+y_2=8\)；由 \(l\) 的斜率为 1：\((y_1-y_2)/(x_1-x_2)=1\)．\(A\)、\(B\) 代入 \(C\) 并相减：\((x_1+x_2)(x_1-x_2)/a^{2}=(y_1+y_2)(y_1-y_2)/b^{2}\)，故 \(b^{2}/a^{2}=(y_1+y_2)(y_1-y_2)/[(x_1+x_2)(x_1-x_2)]=8/2=4\)，即 \(b^{2}=4a^{2}\)，\(e=\sqrt{1+b^{2}/a^{2}}=\sqrt{5}\)．回验：设 \(a^{2}=t\)、\(b^{2}=4t\)，\(y=x+3\) 代入 \(C\)：\(4x^{2}-(x+3)^{2}=4t\)，即 \(3x^{2}-6x-9-4t=0\)，\(\Delta=36+12(9+4t)=144+48t>0\) 恒成立，且 \(x_1+x_2=2\)、\(y_1+y_2=2[(x_1+x_2)/2+3]=8\)，与中点吻合．故选：D．}
\end{ansblock}}

\jiancestem{{\fontsize{11.4pt}{14pt}\selectfont\heihao {\numboldjian 20}．}\hspace{4.1mm}\tieside{中档(知识点二)}\hspace{2.2mm}点\(P(8,1)\)平分双曲线\(x^{2}-4y^{2}=4\)的一条弦，则这条弦所在直线的方程为\kongbai{}．}
\xiexwei{7mm}

{\ansblockgrayfalse
\begin{ansblock}[2章-导-课时19-G4]
% ans:2章-导-课时19-G4
\ansitem{20}{2x−y−15=0}
\ansnote{详解}{（点差法）设弦端点 \(A(x_1,y_1)\)、\(B(x_2,y_2)\)，则 \(x_1+x_2=16\)、\(y_1+y_2=2\)．\(x_1^{2}-4y_1^{2}=4\) 与 \(x_2^{2}-4y_2^{2}=4\) 相减：\((x_1+x_2)(x_1-x_2)-4(y_1+y_2)(y_1-y_2)=0\)，即 \(16(x_1-x_2)=8(y_1-y_2)\)，故 \(k=(y_1-y_2)/(x_1-x_2)=2\)．直线：\(y-1=2(x-8)\)，即 \(2x-y-15=0\)．回验 \(\Delta>0\)：代入得 \(x^{2}-4(2x-15)^{2}=4\)，即 \(15x^{2}-240x+904=0\)，\(\Delta=57600-54240=3360>0\)，弦存在，符合．}
\end{ansblock}}

\jiancestem{{\fontsize{11.4pt}{14pt}\selectfont\heihao {\numboldjian 21}．}\hspace{4.1mm}\tieside{简单(知识点二)}\hspace{2.2mm}已知抛物线方程\(y^{2}=8x\)，求过点\(A(1,1)\)且被该点平分的抛物线的弦所在直线方程．}
\xiexwei{7mm}

{\ansblockgrayfalse
\begin{ansblock}[2章-导-课时19-G5]
% ans:2章-导-课时19-G5
\ansitem{21}{4x−y−3=0}
\ansnote{详解}{（点差法）设弦端点 \(M(x_1,y_1)\)、\(N(x_2,y_2)\)，由中点 \(A(1,1)\)：\(x_1+x_2=2\)、\(y_1+y_2=2\)．若弦垂直于 \(x\) 轴，其中点应在 \(x\) 轴上（纵坐标和为 0），与 \(y_1+y_2=2\) 矛盾，故斜率存在．\(y_1^{2}=8x_1\) 与 \(y_2^{2}=8x_2\) 相减：\((y_1+y_2)(y_1-y_2)=8(x_1-x_2)\)，故 \(k=(y_1-y_2)/(x_1-x_2)=8/(y_1+y_2)=4\)．直线：\(y-1=4(x-1)\)，即 \(4x-y-3=0\)．回验 \(\Delta>0\)：代入得 \((4x-3)^{2}=8x\)，即 \(16x^{2}-32x+9=0\)，\(\Delta=1024-576=448>0\)，弦存在，符合（且 \(x_1+x_2=32/16=2\) 与中点吻合）．}
\end{ansblock}}

\end{ketang}

% ---- 尾块（\tailfill 置 \end{multicols} 前；无参＝内容尾随框，件侧不擅自定高） ----
\tailfill

\end{multicols}
\end{document}
